% ===================================================================
%  ТАБЛИЦЯ № 12 — Вокзал. — Залізниці. — Судна.
%  Traduction 2026 d'apres texto/io, controlee sur texto/fr, texto/en
%  et texto/es. Consigne : voir texto/uk/10-tablycia-01.tex.
%
%  MEME OBSERVATION QU'AU TABLEAU 10, et pour la meme raison : le
%  vocabulaire du chemin de fer est venu tard et par le meme chemin
%  dans les deux langues. « вагон », « перон », « семафор »,
%  « тендер », « стрілка » sont les memes mots. Ce qui differe est ce
%  que le rail a trouve en place : la locomotive est un « паротяг »,
%  la valise une « валіза », le porteur un « носій », la consigne une
%  « камера схову », le garcon de buffet un « кельнер ».
% ===================================================================

%%K t12-apar-1 apar
\VUsaut{4.0mm}
\VUtitre{15.0pt}{60}{ТАБЛИЦЯ № 12}
\VUsaut{2.4mm}
\VUpk{11.4pt}{40}{Вокзал. --- Залізниці. --- Судна.}
\VUsaut{3.4mm}
\textsc{Примітка.} --- \textit{Розклади в кожній країні свої, тому за взірець ми
взяли години} \VUgras{французької таблиці}.

%%K t12-01-1 p
1. --- Я щойно вчинив подорож по Німеччині. Перед від'їздом я купив
\VUgras{розклад} \textsuperscript{(15)}, щоб знати години відходу потягів.
Переконавшись, що найзручніший потяг іде з Парижа о дев'ятій вечора і приходить
до Страсбурга о першій тринадцять, я зібрався в дорогу, а назавтра велів
відвезти себе на вокзал.

%%K t12-02-1 p
2. --- Коли я ввійшов у велику залу відправлення слідом за старим сільським
\VUgras{священником} \textsuperscript{(2)}, який ніс свою \VUgras{дорожню торбу}
\textsuperscript{(3)}, чимало \VUgras{пасажирів} \textsuperscript{(1)} уже
зібралося.

%%K t12-02-2 p
Що мені треба було послати \VUgras{телеграму} \textsuperscript{(5)}, я попросив
\VUgras{контролера} \textsuperscript{(6)} показати мені \VUgras{телеграф}
\textsuperscript{(4)}.

%%K t12-03-1 p
3. --- Перш ніж пройти до \VUgras{зали чекання} \textsuperscript{(7)}, я пішов
узяти зворотний \VUgras{квиток} \textsuperscript{(9)}.

%%K t12-04-1 p
4. --- Коло \VUgras{каси} \textsuperscript{(8)} я чекав недовго, бо народу було
мало. \VUgras{Солдат} \textsuperscript{(10)}, що їхав у відпустку, чемно
поступився мені своєю чергою, тож у мене вистачило часу поставити кілька
запитань \VUgras{начальникові вокзалу} \textsuperscript{(13)}, який розмовляв із
\VUgras{поліційним приставом} \textsuperscript{(12)} і з черговим
\VUgras{жандармом} \textsuperscript{(11)}.

%%K t12-tit-1 sub
\VUpk{10.2pt}{40}{Здача багажу.}
\VUsaut{3.0mm}

%%K t12-05-1 p
5. --- Купивши газету в \VUgras{кіоску} \textsuperscript{(14)}, я велів зважити
свій \VUgras{багаж} \textsuperscript{(17)} --- \VUgras{носій}
\textsuperscript{(16)} щойно привіз його на \VUgras{багажному візку}
\textsuperscript{(19)}. \VUgras{Урядовець} \textsuperscript{(22)} при
\VUgras{вагах} \textsuperscript{(21)} сказав мені, що моя скриня важить тридцять
п'ять кілограмів, тобто на п'ять кілограмів \VUgras{понад норму}, за які
довелося доплатити в \VUgras{багажній касі} \textsuperscript{(23)}.

%%K t12-06-1 p
6. --- Що я приїхав завчасно, у мене був час подивитися на рух пасажирів по
вокзалу. Одні здавали або одержували невеликі торби в \VUgras{камері схову}
\textsuperscript{(25)}; інші справлялися по \VUgras{розкладах}
\textsuperscript{(26)}. Приїжджі поспішали до \VUgras{виходу}
\textsuperscript{(32)} з \VUgras{валізами} \textsuperscript{(24)} в руках. Дехто
чекав коло \VUgras{видачі багажу} \textsuperscript{(27)} і подавав свою
\VUgras{квитанцію} \textsuperscript{(29)}, щоб одержати скрині, \VUgras{ящики}
\textsuperscript{(28)} або \VUgras{кошики} \textsuperscript{(30)}.

%%K t12-07-1 p
7. --- \VUgras{Акцизні урядовці} \textsuperscript{(33)} уважно оглядали пакунки,
чи немає чого до оголошення.

%%K t12-08-1 p
8. --- Інші пасажири прямували до \VUgras{буфету} \textsuperscript{(34)}, де
\VUgras{кельнер} \textsuperscript{(35)} ретельно їх обслуговував. Доводилося
поспішати, бо \VUgras{потяг} \textsuperscript{(36)}, кур'єрський, стояв на
вокзалі недовго.

%%K t12-09-1 p
9. --- Тоді я вийшов на перон відправлення і став шукати прямий \VUgras{вагон}
\textsuperscript{(37)}, щоб не пересідати в дорозі. Я ввійшов у вагон з
коридором, де було \VUgras{купе} \textsuperscript{(38)} для курців і одне,
відведене «для дам, що їдуть самі».

%%K t12-tit-2 sub
\VUpk{10.2pt}{40}{Нічний від'їзд.}
\VUsaut{3.0mm}

%%K t12-10-1 p
10. --- Піднявшись на \VUgras{підніжку} \textsuperscript{(40)}, зачинивши
\VUgras{дверцята} \textsuperscript{(39)}, піднявши \VUgras{штору}
\textsuperscript{(41)} і поклавши свої торби в \VUgras{сітку}
\textsuperscript{(43)}, я сів на \VUgras{диван} \textsuperscript{(42)}.
Побачивши, як \VUgras{ліхтарник} \textsuperscript{(44)} засвічує лампи, я
згадав, що мені ночувати у вагоні, і покликав служника, який дає напрокат
\VUgras{подушки} \textsuperscript{(45)}, щоб узяти одну.

%%K t12-11-1 p
11. --- Кондуктор прийшов перевіряти квитки. Я спитав його, чого ми не рушаємо,
адже час підійшов. Він відповів, що \VUgras{обер-кондуктор}
\textsuperscript{(46)} зайнятий вантаженням велосипедів у \VUgras{багажний
вагон} \textsuperscript{(47)}. До того ж не всі \VUgras{грілки}
\textsuperscript{(48)} ще перемінили.

%%K t12-12-1 p
12. --- Але вирушати нам мало бути скоро, бо \VUgras{кочегар}
\textsuperscript{(50)} на своєму \VUgras{тендері} \textsuperscript{(49)} набирав
вугілля, щоб підтримати вогонь \VUgras{паротяга} \textsuperscript{(54)}, а
\VUgras{машиніст} \textsuperscript{(51)} чекав тільки знаку від
\VUgras{помічника начальника вокзалу} \textsuperscript{(52)}, що стояв на
\VUgras{пероні} \textsuperscript{(53)}.

%%K t12-13-1 p
13. --- \VUgras{Казан} \textsuperscript{(56)} паротяга глухо гудів;
\VUgras{димар} \textsuperscript{(55)} викидав потоки диму навпіл із
\VUgras{парою} \textsuperscript{(60)}. Пролунав пронизливий \VUgras{свисток}
\textsuperscript{(59)}, і \VUgras{поршні} \textsuperscript{(58)} прийшли в рух,
приводячи в хід \VUgras{колеса} \textsuperscript{(57)} важкої машини.

%%K t12-tit-3 sub
\VUsaut{3.0mm}
\VUpk{10.2pt}{40}{Поїхали!}
\VUsaut{3.0mm}

%%K t12-14-1 p
14. --- Потяг, що біг по \VUgras{рейках} \textsuperscript{(64)}, набирав ходу;
\VUgras{перони} \textsuperscript{(65)} скінчилися; ми минули \VUgras{вбиральні}
\textsuperscript{(66)}, а також ряд вагонів, що стояли на іншій \VUgras{колії}
\textsuperscript{(63)}: \VUgras{спальні вагони} \textsuperscript{(67)},
\VUgras{вагон-ресторан} \textsuperscript{(68)}, \VUgras{поштовий вагон}
\textsuperscript{(69)}.

%%K t12-noto-1 noto
\VUnotes{143.2mm}{%
\textit{(*) Примітка.} --- \textit{Подорож, звісно, описано так, як вона
проходила по довоєнному кордоні.}%
}

%%K t12-15-1 p
15. --- Тоді перед нашими очима пройшла \VUgras{товарна станція}
\textsuperscript{(71)}. Під повітками урядовці вантажили \VUgras{товари}
\textsuperscript{(72)}, що відправляються малою швидкістю. \VUgras{Зчіплювачі}
\textsuperscript{(76)} біля \VUgras{водонапірної вежі} \textsuperscript{(73)}
пересували \VUgras{вагони для худоби} \textsuperscript{(75)} і
\VUgras{платформи} \textsuperscript{(79)}, складаючи \VUgras{товарний потяг}
\textsuperscript{(80)}.

%%K t12-16-1 p
16. --- Ми минули останню \VUgras{стрілку} \textsuperscript{(78)}, коло якої
стояв стрілочник; ми пройшли \VUgras{семафор} \textsuperscript{(86)}, і вокзал
зник удалині.

%%K t12-17-1 p
17. --- Незабаром ми виїхали на \VUgras{залізний міст} \textsuperscript{(74)},
який перетинає \VUgras{річку} \textsuperscript{(81)}. Минувши міст, ми побігли
вздовж річки, по якій потужні \VUgras{буксири} \textsuperscript{(82)} тягли
важкі \VUgras{баржі} \textsuperscript{(83)}. Ми бачили й \VUgras{пасажирський
пароплав} \textsuperscript{(84)}, який щойно підходив до \VUgras{пристані}
\textsuperscript{(85)}.

%%K t12-18-1 p
18. --- Промчавши через кілька станцій, ми пройшли кілька \VUgras{тунелів}
\textsuperscript{(87)} і лишили позаду незліченні \VUgras{будиночки}
\textsuperscript{(88)} \VUgras{колійних сторожів}. Заколисаний хитанням потяга,
я заснув міцним сном.

%%K t12-tit-4 sub
\VUsaut{3.0mm}
\VUpk{10.2pt}{40}{Станція Дойч-Аврікур.}
\VUsaut{3.0mm}

%%K t12-19-1 p
19. --- Мене раптом розбудив голос урядовця, який гукав: «Дойч-Аврікур!
\textsuperscript{(*)} Пересадка!» Це був митний огляд і всі докучні прикордонні
формальності. Мені довелося поставити годинник за вокзальним
\VUgras{годинником} \textsuperscript{(89)}, який ішов на п'ятдесят п'ять хвилин
уперед; спросоння я ледве відрізняв \VUgras{велику стрілку}
\textsuperscript{(90)} від \VUgras{маленької} \textsuperscript{(91)}; сам
\VUgras{циферблат} \textsuperscript{(92)} був геть розмитий ранковим туманом.

%%K t12-20-1 p
20. --- Поки митники вели свій огляд, я, позіхаючи, розглядав \VUgras{афіші}
\textsuperscript{(93)}, якими оздоблено стіни вокзалу. Я поснідав, щоб згаяти
час, справився по карті німецької \VUgras{мережі} \textsuperscript{(94)}, чи
далеко ще до Страсбурга, і вернувся до свого вагона, підбадьорений надією скоро
дістатися до мети.
\VUsaut{6.0mm}
\VUfilet{22mm}
